\section{实用扩展指南}

\subsection{从 1 GPU 到 10,000 GPU 的扩展路径}

\begin{figure}[H]
\centering
\includegraphics[width=0.95\textwidth]{slides-images/slide-026.jpg}
\caption{扩展路径：DP $\rightarrow$ FSDP $\rightarrow$ FSDP + 检查点 $\rightarrow$ HSDP $\rightarrow$ 高级并行策略\protect\footnotemark}
\end{figure}
\footnotetext{Slides 第26页。}

\begin{importantbox}{扩展配方}
\begin{enumerate}
    \item \textbf{DP}（$\leq$128 GPU，$\leq$1B 参数）：最简单，将局部 batch size 调到最大化 GPU 内存
    \item \textbf{FSDP}（$>$1B 参数）：权重分片，突破单 GPU 内存限制
    \item \textbf{FSDP + 激活检查点}：激活值内存瓶颈时启用
    \item \textbf{HSDP}（$>$256--512 GPU）：利用网络拓扑，减少跨服务器通信
    \item \textbf{TP + CP + PP}（$>$1,000 GPU，$>$50B 参数，序列 $>$10K）：高级并行策略
\end{enumerate}
\end{importantbox}

\subsection{MFU：分布式训练的北极星指标}

当面对如此多的并行策略和超参数时，\textbf{Model FLOP/s Utilization (MFU)} 是你的指路明灯：

\begin{figure}[H]
\centering
\includegraphics[width=0.95\textwidth]{slides-images/slide-028.jpg}
\caption{MFU 计算方法与基准：Llama 3 在 8K--16K GPU 上达到 38--43\% MFU\protect\footnotemark}
\end{figure}
\footnotetext{Slides 第28页。}

$$\text{MFU} = \frac{\text{模型前向+反向的理论 FLOP 数}}{\text{设备峰值 FLOP/s} \times \text{实际训练迭代时间}}$$

\begin{itemize}
    \item \textbf{$>$30\%}：良好
    \item \textbf{$>$40\%}：优秀（接近当前最优水平）
    \item \textbf{$<$30\%}：可能存在严重瓶颈
\end{itemize}

\begin{warningbox}{更快的 GPU 可能导致更低的 MFU}
从 A100 到 H100，计算吞吐量提升约 3 倍，但内存带宽仅提升约 2 倍。通信速度跟不上计算速度的增长，导致 H100 上的 MFU（$\sim$40\%）反而低于 A100（$\sim$50\%）。这个差距在未来的 GPU 上可能继续扩大。
\end{warningbox}

\subsection{本章小结}

扩展分布式训练是一个渐进的过程，随着 GPU 数量和模型规模的增加逐步引入更复杂的并行策略。MFU 是衡量训练效率的核心指标，所有并行策略的调优都应以最大化 MFU 为目标。

%% ========== 高级并行 ==========

